\begin{quote}
\hfill  ``{\em Everything should be made as simple as possible, but not any simpler}.''

$~$\hfill -- Albert Einstein
\end{quote}

 Closed-Loop Transcription for autonomous, incremental learning and to ensure correctness and self-consistency. This chapter connects to roots of Feedback Control and Game Theory. Content of this chapter might change based on ongoing research...


\begin{definition}[Closed-Loop Representation Learning]\label{def:closed_loop}
    Given data \(\vX\), an \textit{self-consistent representation} is a pair of functions \((f \colon \cX \to \cZ, g \colon \cZ \to \cX)\), such that the \textit{features} \(\vZ = f(\vX)\) are compact and structured, and the \textit{autoencoding features} \(\wh{\vZ} \doteq f \circ g(\vZ)\) is \textit{close} to \(\vZ\). We say that the representation is \textit{distributionally self-consistent} if \(\Law(\vZ) \approx \Law(\wh{\vZ})\); we say that it is \textit{sample-wise} self-consistent if \(\vZ \approx \wh{\vZ}\) with high probability.
\end{definition}


\section{Self-Consistent Representations}
The concept of self-consistent for a system that learns autonomously.

\begin{itemize}
    \item Contrast: supervised representation learning needs human signal (class), self-supervised representation learning needs human signal (masking, augmentation), how to remove human signal from the mixture and get a autonomous system?
    \item For full autonomy, need to work in both \(x\) space (actions, perceptions) and \(z\) space (representations, compression, planning?)
    \item Measure distance in \(x\) space? Deficiencies in \(\ell^{2}\) loss, Wasserstein loss, etc 
    \item Measure distance in \(z\) space? Only possible with compact and structured \(z\) space (say piecewise linear structure, i.e., mixture of subspaces)
    \item Propose self-consistency: representation is consistent wrt itself --- features are consistent wrt their autoencoding 
\end{itemize}

\section{Closed-Loop Transcription}

\subsection{Closed-loop transcription via minimax game}
The closed-loop transcription framework and the game-theoretic formulation: \href{https://www.mdpi.com/1099-4300/24/4/456/htm}{the CTRL work}.

\begin{itemize}
    \item Closed loop transcription framework: feature quality for \(z\) and \(\hat{z}\) plus distance measure for \(z\) and \(\hat{z}\)
    \item \(f\) plays role as encoder, discriminator, sensor; \(g\) plays role as decoder, generator, controller 
    \item \(f\) maximizes above objective, \(g\) minimizes 
    \item optimization strategy: gradient descent-ascent, or two-timescale GDA, or something else?
\end{itemize}

\begin{itemize}
    \item also need to give an intro to game theory somewhere here...
\end{itemize}

\subsection{Theoretical justification for mixed Gaussians}
Theoretical justification for \href{https://arxiv.org/abs/2206.09120}{the case with a mixture of subspaces or Gaussians}.

\section{Applications in Autonomous Continuous Learning}

\subsection{Incremental learning} 
The closed-loop architecture alleviates  catastrophic forgetting, enables \href{https://arxiv.org/abs/2202.05411}{incremental learning}, probably with white-box  backbone networks. 

\subsection{Unsupervised continuous learning}
\href{https://arxiv.org/abs/2210.16782}{unsupervised learning}, probably with white-box backbone networks.

\subsection{Other applications} 
to be determined...

